\chapter{Francia: sabotaje de la refrigeración comercial}
\label{chap:incidente4}

\section{Contexto y justificación de selección}
Durante la noche del 11 al 12 de octubre de 2026,
varios cientos de supermercados franceses sufrieron
un ataque contra sus sistemas de refrigeración.
La activación de ciclos de descongelación alteró
la conservación de los productos y provocó pérdidas
de mercancía~\cite{onyphe2026francia}.
Se selecciona por mostrar cómo un incidente de
ciberseguridad puede afectar a procesos físicos,
a la distribución de alimentos y a la actividad
cotidiana de comercios y consumidores.

\section{Desarrollo del incidente y CVE/CVSS}
El CERT-FR confirmó compromisos de controladores
Danfoss AK-SM y AK-SC, además de sistemas Copeland
y Dixell XWEB~\cite{certfr2026frigorificos}.
Estos equipos permiten supervisar instalaciones y
modificar parámetros operativos desde una interfaz
web. En condiciones normales facilitan el mantenimiento
remoto; su compromiso permite que una acción digital
altere el funcionamiento de la
refrigeración~\cite{claroty2026danfoss}.

Existían vulnerabilidades documentadas en algunos
modelos. Claroty describió en el AK-SM 800A un mecanismo
alternativo de autenticación basado en un código
derivado de la fecha, así como fallos posteriores a
la autenticación que permitían ejecutar comandos
o modificar la configuración del servidor
web~\cite{claroty2026danfoss}.

Como antecedente técnico, \textbf{CVE-2025-41450}
afecta a la autenticación de la serie AK-SM 8xxA
anterior a la versión 4.2. Danfoss publica
\textbf{CVSS v3.1: 8,2, alto}, con el
vector~\cite{danfoss2025cve41450}:\par
\noindent{\small\ttfamily
    CVSS:3.1/AV:N/AC:H/PR:N/UI:R/S:C/C:H/I:L/A:H}\par
\noindent Esta puntuación corresponde al fallo conocido,
no al conjunto del incidente. Las fuentes consultadas
no permiten demostrar que fuera la vía utilizada
en todos los equipos atacados. Sin identificar
la vulnerabilidad concreta, no resulta defendible
asignar un CVSS propio al sabotaje.

\section{Impacto y mitigación}
La pérdida de refrigeración obligó a retirar alimentos,
con perjuicios económicos para los establecimientos
y desperdicio de productos~\cite{onyphe2026francia}.
Para los consumidores, la alteración de la cadena
de frío representa un riesgo sanitario si los productos
afectados llegan a consumirse; no se afirma que hubiera
intoxicaciones confirmadas.
El impacto combina manipulación de parámetros,
pérdida del servicio y consecuencias materiales.

El CERT-FR recomendó eliminar la exposición directa
a Internet, restablecer el acceso mediante una VPN
segura o una red protegida, actualizar el firmware
y cambiar las contraseñas predeterminadas.
Para equipos Danfoss compatibles, indicó utilizar
el nivel de seguridad \textit{Strict}.
Advirtió también que restaurar un equipo sin corregir
su exposición o vulnerabilidades no evita una nueva
intrusión, especialmente si ha terminado su
soporte~\cite{certfr2026frigorificos}.
Como prevención adicional, se propone supervisión
independiente de temperatura y alarmas ante cambios
anómalos de configuración.

\section{Análisis crítico}
La comodidad del mantenimiento remoto debe evaluarse
junto con las consecuencias de perder el control
del proceso físico. La protección requiere coordinar
al comercio, al mantenedor y al fabricante, definiendo
responsabilidades sobre accesos, actualizaciones
y respuesta.
Además, encontrar dispositivos expuestos o fallos
compatibles no demuestra el mecanismo exacto del
ataque: separar esas evidencias evita convertir
una hipótesis técnica en una causa confirmada.